% Appendix draft for the main branch. Insert before \bibliography{...}.
% Model architectures and capacity counts are included from the code-audited
% English fragment below.

\section{完整实验细节}
\label{app:experimental_details}

\subsection{数据预处理与联邦划分}

本文使用 CIFAR-10、CIFAR-100 和 Tiny-ImageNet 三个图像分类数据集。
所有图像先转换为张量，并对三个 RGB 通道统一采用均值
$(0.5,0.5,0.5)$ 和标准差 $(0.5,0.5,0.5)$ 进行归一化。
训练过程中不使用随机裁剪、随机翻转或其他数据增强。
CIFAR-10/100 的官方训练集与测试集首先合并，然后执行联邦客户端划分；
Tiny-ImageNet 仅使用官方训练目录中的 $100{,}000$ 张图像，未使用官方验证集。
完成客户端划分后，每个客户端的数据独立地按 $75\%:25\%$ 随机划分为本地训练集和本地测试集。
因此，本文的本地测试集并非沿用 CIFAR 的官方测试集，而是由合并后的样本在客户端内部重新划分得到。

\begin{table*}[t]
\centering
\caption{数据预处理与本地训练/测试划分。}
\label{tab:appendix_data_protocol}
\small
\setlength{\tabcolsep}{6pt}
\renewcommand{\arraystretch}{1.12}
\begin{tabularx}{\textwidth}{@{}lcccccX@{}}
\toprule
\textbf{数据集} & \textbf{类别数} & \textbf{输入尺寸} & \textbf{划分前样本数} & \textbf{本地训练比例} & \textbf{本地测试比例} & \textbf{划分前使用的数据池} \\
\midrule
CIFAR-10 & 10 & $3\times32\times32$ & 60,000 & 75\% & 25\% & 官方训练集与测试集合并 \\
CIFAR-100 & 100 & $3\times32\times32$ & 60,000 & 75\% & 25\% & 官方训练集与测试集合并 \\
Tiny-ImageNet & 200 & $3\times64\times64$ & 100,000 & 75\% & 25\% & 仅使用官方训练目录 \\
\bottomrule
\end{tabularx}
\end{table*}

本文考虑 Dirichlet 和病理标签划分两类统计异构。
在 Dirichlet 设置中，对每个类别从参数为 $\alpha$ 的 Dirichlet 分布中采样其在客户端之间的分配比例，
并设置 $\alpha\in\{0.3,0.5,0.8,1.0\}$。
划分程序要求每个客户端至少具有足以产生一个测试批次的样本。
在病理划分中，每个客户端只观察固定数量的类别，并采用均衡样本分配；
不同数据集使用的每客户端类别数见表~\ref{tab:appendix_partition}。
数据集生成阶段固定 Python 和 NumPy 随机种子为 0，且生成后的划分以独立的
\texttt{npz} 文件持久化，所有方法复用完全相同的文件。

\begin{table}[t]
\centering
\caption{统计异构设置。$q$ 表示病理划分中每个客户端包含的类别数。}
\label{tab:appendix_partition}
\small
\setlength{\tabcolsep}{5pt}
\renewcommand{\arraystretch}{1.10}
\begin{tabular}{lcc}
\toprule
\textbf{数据集} & \textbf{Dirichlet 系数 $\alpha$} & \textbf{病理划分 $q$} \\
\midrule
CIFAR-10 & $\{0.3,0.5,0.8,1.0\}$ & $\{2,4,6,8\}$ \\
CIFAR-100 & $\{0.3,0.5,0.8,1.0\}$ & $\{20,40,60,80\}$ \\
Tiny-ImageNet & $\{0.3,0.5,0.8,1.0\}$ & $\{40,80,120,160\}$ \\
\bottomrule
\end{tabular}
\end{table}

对应的数据生成命令分别遵循以下形式：
\begin{verbatim}
python generate_<dataset>.py --niid 1 --balance 0 --partition dir \
    --alpha <alpha>
python generate_<dataset>.py --niid 1 --balance 1 --partition pat \
    --cpc <classes_per_client>
\end{verbatim}
需要注意，数据生成脚本使用参数名 \texttt{--alpha}，训练脚本使用
\texttt{--dir\_alpha} 读取对应的预生成目录。

\subsection{Model Architectures and Capacity Configurations}
\label{app:model_capacity}

\paragraph{Base model architectures.}
The experiments use two backbone families.  The CNN is used in the main
experiments on CIFAR-10, CIFAR-100, and Tiny-ImageNet, whereas ResNet-18 is
used to evaluate whether FedLAP extends to a deeper backbone.  Within each
FedLAP low-rank family, all clients retain the same macro-architecture and
differ only in the realized ranks of the factorized layers.  CIFAR-10 and
CIFAR-100 use $32\times32$ inputs.  Tiny-ImageNet uses $64\times64$ inputs;
the CNN therefore changes the input dimension of its first fully connected
layer from 800 to 5,408, and ResNet-18 uses a larger stem.  The classifier
output dimensions are 10, 100, and 200 for CIFAR-10, CIFAR-100, and
Tiny-ImageNet, respectively.  Table~\ref{tab:appendix_model_architectures}
reports the instantiated architectures.

\begin{table*}[t]
\centering
\caption{FedLAP model architectures used in the experiments.  ``Factorized''
denotes trainable low-rank parameterization; all other trainable layers retain
their full parameterization.}
\label{tab:appendix_model_architectures}
\scriptsize
\setlength{\tabcolsep}{4pt}
\renewcommand{\arraystretch}{1.12}
\begin{tabularx}{\textwidth}{@{}lXXp{2.7cm}X@{}}
\toprule
\textbf{Layer or stage} & \textbf{CIFAR-10/100 configuration} &
\textbf{Tiny-ImageNet configuration} & \textbf{Parameterization} &
\textbf{Normalization or semantic-alignment role} \\
\midrule
\multicolumn{5}{c}{\textbf{Part A: CNN}} \\
\midrule
Input & $3\times32\times32$ & $3\times64\times64$ & -- & -- \\
Conv1 & $5\times5$, $3\rightarrow16$, stride 1 & Same & Full & No normalization; followed by ReLU \\
Pool1 & $2\times2$ max pooling, stride 2 & Same & -- & -- \\
Conv2 & $5\times5$, $16\rightarrow32$, stride 1 & Same & Factorized & No normalization; followed by ReLU \\
Pool2 & $2\times2$ max pooling, stride 2 & Same & -- & -- \\
Flatten & $32\times5\times5=800$ & $32\times13\times13=5{,}408$ & -- & -- \\
FC1 & $800\rightarrow2{,}000$ & $5{,}408\rightarrow2{,}000$ & Factorized & Followed by ReLU \\
FC2 & $2{,}000\rightarrow512$ & Same & Factorized & Produces the 512-dimensional image representation \\
Classifier & $512\rightarrow C$, $C\in\{10,100\}$ & $512\rightarrow200$ & Full & ReLU is applied before the classifier \\
CLIP semantic alignment & Directly aligns the 512-dimensional FC2 output & Same & No adapter & MSE to the frozen 512-dimensional class text anchor \\
\midrule
\multicolumn{5}{c}{\textbf{Part B: ResNet-18}} \\
\midrule
Input & $3\times32\times32$ & $3\times64\times64$ & -- & -- \\
Stem & $3\times3$, $3\rightarrow64$, stride 1; no max pool & $7\times7$, $3\rightarrow64$, stride 2; $3\times3$ max pool, stride 2 & Full & LayerNorm2d and ReLU after the stem convolution \\
Stage 1 & $[3\times3,64;\ 3\times3,64]\times2$, stride 1 & Same & Both main-branch convolutions in every block are factorized & LayerNorm2d after each convolution \\
Stage 2 & $[3\times3,128;\ 3\times3,128]\times2$, first block stride 2 & Same & Both main-branch convolutions in every block are factorized & LayerNorm2d after each convolution \\
Stage 3 & $[3\times3,256;\ 3\times3,256]\times2$, first block stride 2 & Same & Both main-branch convolutions in every block are factorized & LayerNorm2d after each convolution \\
Stage 4 & $[3\times3,512;\ 3\times3,512]\times2$, first block stride 2 & Same & Both main-branch convolutions in every block are factorized & LayerNorm2d after each convolution \\
Downsampling shortcuts & $1\times1$ convolutions in the first blocks of Stages 2--4 & Same & Full & Followed by LayerNorm2d \\
Global pooling & Adaptive average pooling to $1\times1$ and flattening & Same & -- & Produces a 512-dimensional final representation \\
Classifier & $512\rightarrow C$, $C\in\{10,100\}$ & $512\rightarrow200$ & Full & No factorization \\
CLIP semantic alignment & Pooled outputs of Stages 1--4 have dimensions $64,128,256,512$ and are mapped to 512 dimensions & Same & Four full linear maps: $64/128/256/512\rightarrow512$ & Client-local maps to four frozen text anchors; not uploaded or aggregated \\
\bottomrule
\end{tabularx}
\end{table*}

\paragraph{Low-rank model construction.}
FedLAP constructs each capacity family from one common macro-architecture.
For the CNN, the factorized-layer set comprises Conv2, FC1, and FC2; Conv1,
the classifier, and every bias remain full.  For ResNet-18, all 16
main-branch $3\times3$ convolutions in the eight residual blocks are
factorized.  The stem, the three $1\times1$ downsampling shortcuts, the
classifier, all normalization parameters, and all biases remain full.  A
convolutional kernel is unfolded in the implementation from
$[C_{\mathrm{out}},C_{\mathrm{in}},K,K]$ to a
$(C_{\mathrm{out}}K)\times(C_{\mathrm{in}}K)$ matrix before its retained rank
is selected.  The same configured rank ratio is used by all factorized layers
of a client, while integer ranks are obtained by rounding that ratio times the
maximum unfolded rank and enforcing a minimum rank of one.  The four ResNet
semantic maps are client-local auxiliary modules rather than members of the
low-rank model: together they contain 493,568 trainable parameters, are not
factorized, and are excluded from the model-capacity counts in
Table~\ref{tab:appendix_capacity_configurations}.  The CNN requires no
semantic adapter.  After server aggregation, the reconstructed model is
factorized again at the destination client's configured capacity using the
same recovery and decomposition routines as during initialization.  This is
the experimental realization of the low-rank construction defined in the
main paper; its decomposition equations are not repeated here.

\paragraph{Capacity configurations.}
Table~\ref{tab:appendix_capacity_configurations} reports parameter counts
obtained by instantiating the current code and summing tensor elements in
\texttt{model.parameters()}; every listed model parameter is trainable.  The
counts include full biases, normalization parameters, and classifiers, but
exclude the frozen CLIP text encoder, cached text anchors, and the client-local
ResNet semantic maps described above.  Since the classifier size depends on
the number of classes, the CIFAR column reports CIFAR-10 and CIFAR-100 counts
separately as ``C10/C100.''  Capacity levels are ordered from the smallest
($L1$) to the largest ($L5$).

\begin{table*}[t]
\centering
\caption{Five client-capacity configurations.  Parameter counts are in
millions (M); paired CIFAR entries are reported as CIFAR-10/CIFAR-100.}
\label{tab:appendix_capacity_configurations}
\scriptsize
\setlength{\tabcolsep}{5pt}
\renewcommand{\arraystretch}{1.12}
\begin{tabularx}{\textwidth}{@{}llccccc@{}}
\toprule
\textbf{Capacity} & \textbf{Rank ratio} &
\textbf{FedLAP low-rank model, C10/C100 (M)} &
\textbf{FedLAP low-rank model, Tiny (M)} &
\textbf{Corresponding full-rank model, C10/C100 (M)} &
\textbf{Corresponding full-rank model, Tiny (M)} &
\textbf{Clients} \\
\midrule
\multicolumn{7}{c}{\textbf{Part A: CNN}} \\
\midrule
L1 & 0.15 & 0.541 / 0.587 & 2.525 & 0.676 / 0.722 & 3.078 & 4 \\
L2 & 0.25 & 0.895 / 0.941 & 4.137 & 1.070 / 1.116 & 4.854 & 4 \\
L3 & 0.35 & 1.249 / 1.295 & 5.748 & 1.333 / 1.379 & 6.038 & 4 \\
L4 & 0.37 & 1.320 / 1.366 & 6.070 & 1.839 / 1.885 & 6.545 & 4 \\
L5 & 0.90 & 3.200 / 3.246 & 14.616 & 3.302 / 3.348 & 14.920 & 4 \\
\midrule
\multicolumn{7}{c}{\textbf{Part B: ResNet-18}} \\
\midrule
L1 & 0.12 & 2.728 / 2.774 & 2.833 & 2.932 / 2.978 & 3.033 & 4 \\
L2 & 0.20 & 4.429 / 4.475 & 4.534 & 4.535 / 4.581 & 4.637 & 4 \\
L3 & 0.29 & 6.333 / 6.379 & 6.438 & 6.487 / 6.533 & 6.590 & 4 \\
L4 & 0.40 & 8.664 / 8.710 & 8.769 & 8.787 / 8.834 & 8.892 & 4 \\
L5 & 0.50 & 10.787 / 10.833 & 10.892 & 11.174 / 11.220 & 11.279 & 4 \\
\bottomrule
\end{tabularx}
\begin{minipage}{\textwidth}
\footnotesize
\textit{Note.} The corresponding full-rank CNN models for $L1$--$L5$ are
\texttt{CNN\_5\_512} through \texttt{CNN\_1\_512}, with the
\texttt{\_tiny} variants used for Tiny-ImageNet.  The corresponding full-rank
ResNet-18 models use base widths $32,40,48,56,64$, respectively, and a
512-dimensional output feature.  These model families have approximately
matched parameter counts rather than strictly identical parameter counts.
\end{minipage}
\end{table*}

\paragraph{Client capacity assignment.}
The implementation selects the client model using
\texttt{args.models[cid \% len(args.models)]}.  With 20 clients and five
models, this deterministically assigns four clients to each level:
$L1=\{4,9,14,19\}$, $L2=\{3,8,13,18\}$,
$L3=\{2,7,12,17\}$, $L4=\{1,6,11,16\}$, and
$L5=\{0,5,10,15\}$.  The assignment is fixed by client identifier and is
therefore unchanged across random seeds and data partitions.  It is also
independent of each client's sample count and label distribution.  The
five-model comparison families are listed in the same largest-to-smallest
order in \texttt{main.py}, so their parameter scales correspond by client
identifier to the FedLAP low-rank levels, subject to the approximate rather
than exact matches reported in Table~\ref{tab:appendix_capacity_configurations}.



\subsection{统一训练协议}

除特别说明外，所有方法采用表~\ref{tab:appendix_training_protocol} 中的统一协议。
每轮随机选择规定比例的客户端，客户端在本地执行 5 个 epoch 的训练后上传模型。
客户端本地优化器为不带动量和权重衰减的 SGD，基础学习率为 0.005。
训练数据加载器启用随机打乱，测试数据加载器不打乱，且均不丢弃最后一个不完整批次。
主实验训练 100 个通信轮次；客户端参与率实验在
$\{0.2,0.4,0.8\}$ 上训练 500 个通信轮次，其余设置不变。
模型在每个通信轮次进行一次评估，不采用早停。

\begin{table}[t]
\centering
\caption{所有方法共享的训练设置。}
\label{tab:appendix_training_protocol}
\small
\setlength{\tabcolsep}{6pt}
\renewcommand{\arraystretch}{1.10}
\begin{tabular}{lc}
\toprule
\textbf{设置} & \textbf{取值} \\
\midrule
客户端总数 & 20 \\
默认参与率 & 1.0 \\
参与率实验 & $\{0.2,0.4,0.8\}$ \\
主实验通信轮数 & 100 \\
参与率实验通信轮数 & 500 \\
本地 epoch 数 & 5 \\
本地 batch size & 16 \\
客户端优化器 & SGD \\
基础学习率 & 0.005 \\
SGD momentum / weight decay & 0 / 0 \\
评估间隔 & 每 1 个通信轮次 \\
客户端掉线率 / 慢客户端比例 & 0 / 0 \\
\bottomrule
\end{tabular}
\end{table}

\subsection{FedLAP 本地目标与语义锚点}

对于客户端 $i$ 的样本 $(x,y)$，FedLAP 的本地目标为
\begin{equation}
\mathcal{L}_i=
\mathcal{L}_{\mathrm{CE}}(f_i(x),y)
+\lambda_a\mathcal{L}_{\mathrm{anchor}}
+\lambda_r\sum_{l\in\mathcal{D}}\lVert U_{i,l}V_{i,l}\rVert_F^2,
\end{equation}
其中 $\mathcal{D}$ 表示低秩参数化层，$\lambda_a=1$，CNN 实验使用
$\lambda_r=10^{-3}$。反向传播后对全部参与训练的参数执行最大范数为 10 的梯度裁剪。

语义锚点由冻结的 CLIP ViT-B/32 文本编码器产生，提示模板固定为
\texttt{a photo of \{class name\}}。本文不使用 CLIP 图像编码器。
文本锚点在首次生成后缓存到本地，后续实验直接复用，因此其计算不计入每轮联邦训练开销。
对于 CNN，使用最终的 512 维图像特征与类别文本锚点计算均方误差：
\begin{equation}
\mathcal{L}_{\mathrm{anchor}}^{\mathrm{CNN}}
=\lVert h_i(x)-e_y\rVert_2^2.
\end{equation}
对于 ResNet-18，从四个残差 stage 的输出分别进行全局平均池化，
并通过四个客户端本地的线性映射层投影到 CLIP 文本空间。
四个文本锚点分别取自 CLIP 文本 Transformer 的
$1/4$、$2/4$、$3/4$ 和 $4/4$ 深度，四项 MSE 损失等权平均。
这些映射层使用基础学习率 0.005，只在对应客户端本地更新，不上传服务器，也不参与跨客户端聚合。

FedLAP 表示学习率匹配版本，所有低秩因子和其他模型参数均使用 0.005；
FedLAP\textsuperscript{*} 表示非对称学习率版本，其中
$\eta_V=0.005$、$\eta_U=0.1\eta_V=0.0005$，其余模型参数仍使用 0.005。

\subsection{容量感知有序前缀训练}

低秩层在秩维度上采用嵌套前缀掩码，仅激活前 $r$ 个秩分量。
默认动态容量感知策略由训练进度 $p=t/T$ 控制：当 $p\leq0.3$ 时使用客户端完整秩；
当 $0.3<p<0.8$ 时，以 $(p-0.3)/(0.8-0.3)$ 的概率启用容量感知采样；
当 $p\geq0.8$ 时始终启用容量感知采样。
CNN 的候选秩率和采样概率见表~\ref{tab:appendix_cnn_dropout}。
所有候选秩均不超过客户端自身容量，因此较小客户端不会采样无法容纳的更大秩。

\begin{table*}[t]
\centering
\caption{CNN 容量感知前缀采样。每行括号内为候选秩率及对应概率。}
\label{tab:appendix_cnn_dropout}
\small
\setlength{\tabcolsep}{6pt}
\renewcommand{\arraystretch}{1.10}
\begin{tabularx}{\textwidth}{@{}cXX@{}}
\toprule
\textbf{客户端秩率} & \textbf{候选秩率} & \textbf{采样概率} \\
\midrule
0.90 & $(0.90,0.37,0.35,0.25,0.15)$ & $(0.65,0.15,0.10,0.07,0.03)$ \\
0.37 & $(0.37,0.35,0.25,0.15)$ & $(0.65,0.15,0.12,0.08)$ \\
0.35 & $(0.35,0.25,0.15)$ & $(0.70,0.20,0.10)$ \\
0.25 & $(0.25,0.15)$ & $(0.75,0.25)$ \\
0.15 & $(0.15)$ & $(1.00)$ \\
\bottomrule
\end{tabularx}
\end{table*}

\subsection{服务器个性化聚合细节}

服务器按客户端本地训练样本数计算全局可靠性权重
$\alpha_j=n_j/\sum_{k\in\mathcal{S}^t}n_k$。
对客户端上传后的低秩参数与该客户端本轮实际训练起点做差，得到逐层低秩更新。
对于可低秩化层，在 $V$ 因子的更新上计算客户端相似度；
当两个客户端秩不同，只比较二者共同拥有的有序秩前缀。
记第 $l$ 个逻辑层的共同前缀更新为 $\Delta V_{i,l}^{\cap}$，则
\begin{equation}
s_{ij,l}=\frac{\langle\Delta V_{i,l}^{\cap},\Delta V_{j,l}^{\cap}\rangle}
{\lVert\Delta V_{i,l}^{\cap}\rVert_2\lVert\Delta V_{j,l}^{\cap}\rVert_2+\epsilon},
\qquad
p_{ij,l}=\frac{\exp(s_{ij,l}/\tau)}{\sum_{k\in\mathcal{S}^t}\exp(s_{ik,l}/\tau)}.
\end{equation}
当前实现不在个性化分支中额外乘样本量缩放，也不增加目标客户端自保留偏置。

第 $l$ 个逻辑层的个性化比例随深度线性增加：
\begin{equation}
d_l=d_{\max}\frac{l}{L},\qquad d_{\max}=0.7,
\end{equation}
最终贡献权重为
\begin{equation}
\omega_{ij,l}=(1-d_l)\alpha_j+d_lp_{ij,l}.
\end{equation}
服务器将各客户端低秩模型恢复到满秩有效权重空间，使用 $\omega_{ij,l}$ 对同一逻辑层的全部参数进行加权，
再将所得个性化满秩模型按目标客户端秩率重新分解后下发。
CNN 中每个卷积层或全连接层构成一个逻辑层；ResNet-18 使用 18 个逻辑聚合单元。
首层与分类器等不存在低秩 $V$ 的层使用满秩权重更新计算相似度。
聚合温度的默认设置为 $\tau=1$。

\subsection{基线超参数}

所有基线均使用相同的数据文件、客户端容量分配和公共本地训练协议。
方法特有设置见表~\ref{tab:appendix_baseline_hparams}；未列出的参数沿用对应实现默认值。
其中服务器端额外优化过程属于方法自身机制，不替换客户端统一使用的 SGD 优化器。

\begin{table*}[t]
\centering
\caption{基线方法的主要专属超参数。}
\label{tab:appendix_baseline_hparams}
\small
\setlength{\tabcolsep}{5pt}
\renewcommand{\arraystretch}{1.12}
\begin{tabularx}{\textwidth}{@{}lX@{}}
\toprule
\textbf{方法} & \textbf{专属设置} \\
\midrule
FedGH & 服务器分类头学习率 0.005；服务器训练 5 epoch \\
LG-FedAvg & 无额外调节参数 \\
FedTGP & 原型损失系数 10；服务器训练 20 epoch；margin threshold 100；特征维度 512 \\
FedProto & 原型正则系数 10 \\
FML & $\alpha=0.5$，$\beta=0.5$ \\
FD & 蒸馏系数 1 \\
FedKD & mentee 学习率 0.005；$T_{\mathrm{start}}=0.95$；$T_{\mathrm{end}}=0.98$；特征维度 512 \\
FedGen & 噪声维度 32；生成器学习率 0.05；隐藏维度 512；服务器训练 20 epoch；特征维度 512 \\
FedMRL & 特征维度 512；子特征维度 128 \\
pFedAFM & 融合系数优化学习率 0.01 \\
FedSPU & 沿用方法实现的默认掩码与聚合设置 \\
FedRE & 服务器头学习率 0.01；服务器训练 1 epoch；服务器头 batch size 10；每客户端上传 1 个纠缠表示 \\
\bottomrule
\end{tabularx}
\end{table*}

\subsection{评价、结果保存与运行环境}

客户端 $i$ 使用其个性化模型在本地测试集上计算正确分类数 $c_i$ 和测试样本数 $m_i$。
表格中的测试准确率实际采用按测试样本数加权的微平均：
\begin{equation}
\mathrm{Acc}=\frac{\sum_{i=1}^{N}c_i}{\sum_{i=1}^{N}m_i}\times100\%.
\end{equation}
代码同时计算各客户端准确率的标准差用于诊断，但主表报告上述微平均准确率。
每轮测试准确率、训练损失和通信开销保存为 HDF5/JSON 日志，训练结束后额外保存服务器模型和全部客户端个性化模型。

\begin{table}[t]
\centering
\caption{软硬件运行环境。提交前请使用实际服务器信息替换方括号内容。}
\label{tab:appendix_environment}
\small
\setlength{\tabcolsep}{6pt}
\renewcommand{\arraystretch}{1.10}
\begin{tabular}{ll}
\toprule
\textbf{项目} & \textbf{版本或配置} \\
\midrule
操作系统 & \texttt{[待填写，例如 Ubuntu 22.04]} \\
GPU & \texttt{[待填写 GPU 型号与显存]} \\
CPU / 内存 & \texttt{[待填写]} \\
Python & \texttt{[待填写]} \\
PyTorch / torchvision & \texttt{[待填写]} \\
CUDA / cuDNN & \texttt{[待填写]} \\
CLIP 实现 & OpenAI CLIP，ViT-B/32 \\
\bottomrule
\end{tabular}
\end{table}

% Reproducibility checks before using this draft in the final manuscript:
% 1. Current main defaults dynamic_capacity for both CNN and ResNet-18, while
%    the manuscript currently says ResNet-18 uses original ordered dropout.
% 2. Some recorded ResNet commands use regular_lamda=1e-4 and aggregate_tau=5,
%    whereas the main text states 1e-3 and a shared setting. Confirm from the
%    JSON/H5 args of the runs used in the final tables.
% 3. main.py currently calls torch.manual_seed(0), but the full set_seed(0)
%    helper (Python/NumPy/PyTorch) is commented out and cuDNN is nondeterministic.
% 4. serverCLIP.py iterates range(global_rounds + 1); verify the intended count
%    of local-train/aggregate steps when describing exactly 100/500 rounds.
